\documentclass[10pt]{article}
\usepackage[margin=0.8in]{geometry}
\usepackage{amsmath,amssymb,booktabs,array,url}
\usepackage[hidelinks]{hyperref}
\newcommand{\confirm}[1]{\par\smallskip\noindent\fbox{\parbox{0.97\linewidth}{\textbf{Please confirm:} #1}}\par\medskip}
\title{Advisor verification packet:\\ equations and statistics of the CBF-QP filter study}
\author{Al Mahdi Sallam (prepared for Dr.~Mohandes and Dr.~Brahmi)}
\date{Draft, 2 November 2026. Repository \texttt{mahdisallam860/cbf-marl-filter-study}, branch \texttt{claude/elegant-hopper-v8jxlv}}
\begin{document}
\maketitle
\noindent This packet lists the six items whose correctness the manuscript relies on. Each gives the
statement, a derivation, the code that implements it and what to confirm. Item numbers match the
sign-off table at the end. The final test was run on 25--26 October 2026 under a protocol the author
locked without advisor review (commit \texttt{9baf6d7}); nothing below changes a reported number.

\section*{Setting}
Unicycle robots, pose $(x,y,\theta)$, command $u=(v,\omega)\in\mathcal U=[-0.5,0.5]\,\text{m/s}\times[-1,1]\,\text{rad/s}$,
held for $\Delta=0.1$\,s. Offset point $p=(x+\ell\cos\theta,\,y+\ell\sin\theta)$, $\ell=0.05$\,m, so
$\dot p=J(\theta)u$ with $J=\begin{bmatrix}\cos\theta&-\ell\sin\theta\\ \sin\theta&\ell\cos\theta\end{bmatrix}$, $\det J=\ell$.
Barriers $h_o=\|p-c\|^2-D_o^2$, $D_o=R+\ell+R_o=0.35$\,m; $h_{ij}=\|p_i-p_j\|^2-D_{rr}^2$, $D_{rr}=2(R+\ell)=0.30$\,m.
Since $\|x-p\|=\ell$, $h\ge0$ implies the body discs (radius $R=0.10$\,m) do not overlap.

\section*{C1. Lemma 1 (one-sided sampled-data condition)}
\textbf{Statement.} Let $u$ be constant on $[t_k,t_k+\Delta]$, $\ddot h\ge-M$ there ($M\ge0$),
$\alpha\Delta\le1$, $h(t_k)\ge0$. If $\dot h(t_k)+\alpha h(t_k)\ge\delta:=M\Delta/2$, then
$h(t_k+s)\ge(1-\alpha s)h(t_k)\ge0$ for all $s\in[0,\Delta]$.

\textbf{Proof.} By Taylor's theorem with the lower bound on $\ddot h$,
$h(t_k+s)\ge h+\dot h\,s-\tfrac M2 s^2$ (values at $t_k$). Using $\dot h\ge-\alpha h+M\Delta/2$,
\[
h(t_k+s)\ \ge\ (1-\alpha s)h+\tfrac{M s}{2}(\Delta-s)\ \ge\ (1-\alpha s)h\ \ge\ 0,
\]
since $s\le\Delta$ and $\alpha s\le\alpha\Delta\le1$. $\square$
Induction over periods gives $h\ge0$ for all $t$ provided every QP is feasible (assumption A9) and
$h(0)\ge0$ (A5). This is Theorem~3 of Breeden, Garg and Panagou, IEEE L-CSS 6:367--372 (2022), there
with safe set $\{h\le0\}$ and $\gamma\in(0,1]$, which corresponds to $\alpha\Delta\le1$.
At $\alpha\Delta=1$ the endpoint bound is $h(t_k+\Delta)\ge0$, not strictly positive.

\noindent Code: \texttt{cbf\_study/cbf\_qp.py} (module docstring, \texttt{FilterConfig.validate} enforces $\alpha\Delta\le1$).
\confirm{the Taylor step, the use of a one-sided bound, the induction under A5 and A9, and the mapping to Breeden et al.\ Thm.~3.}

\section*{C2. The margins $\delta_o(d)$ and $\delta_{ij}(d)$ (eq.~4)}
Under a held command, $\|\dot p\|^2=v^2+\ell^2\omega^2\le V_p^2$ with $V_p=\sqrt{0.5^2+0.05^2}\approx0.5025$\,m/s.
Since $\dot\theta=\omega$ is constant, $\ddot p=\omega\,J'(\theta)u$ and $\|J'(\theta)u\|=\|\dot p\|$,
so $\|\ddot p\|\le\bar\omega V_p=:A_p$. Within the period $\|p(t)-c\|\le d+V_p\Delta$, with $d=\|p(t_k)-c\|$.
\[
\ddot h_o=2\|\dot p\|^2+2(p-c)^\top\ddot p\ \ \Rightarrow\ \ |\ddot h_o|\le 2V_p^2+2(d+V_p\Delta)A_p=:M_o(d).
\]
For a pair, relative speed $\le2V_p$, relative acceleration $\le2A_p$, $\|p_i-p_j\|\le d+2V_p\Delta$:
$|\ddot h_{ij}|\le 8V_p^2+4(d+2V_p\Delta)A_p=:M_{ij}(d)$. The implementation uses
$\delta=M\Delta/2$ with these two-sided $M$; Lemma~1 needs only $\ddot h\ge-M$, for which
$2(d+V_p\Delta)A_p$ (obstacles) suffices, so the $2V_p^2$ and $8V_p^2$ terms are extra conservatism
(about $13\times$ the tight one-sided obstacle margin at $d=D_o$; \texttt{results/design/margin\_constants.json}).
The row is $\dot h_o=g_o^\top u$ with $g_o=2J^\top(p-c)$, i.e.\ $g_o^\top u\ge-\alpha h_o+\delta_o(d)$.

\noindent Code: \texttt{cbf\_qp.margin\_obstacle}, \texttt{cbf\_qp.margin\_robot}, \texttt{bmatch.barrier\_rows}.
\confirm{$\|\ddot p\|\le A_p$ under a held command; both bounds on $|\ddot h|$; that two-sided bounds are valid (if conservative) for the lemma.}

\section*{C3. Equal split between robots}
Robot $i$ imposes $g_{ij}^\top u_i\ge-\tfrac\alpha2 h_{ij}+\tfrac12\delta_{ij}(d)$ with $g_{ij}=2J_i^\top(p_i-p_j)$;
robot $j$ imposes the same with $i\leftrightarrow j$. Since $h_{ij}=h_{ji}$, $d$ is symmetric, and
$\dot h_{ij}=2(p_i-p_j)^\top(\dot p_i-\dot p_j)=g_{ij}^\top u_i+g_{ji}^\top u_j$, adding the two rows gives
$\dot h_{ij}\ge-\alpha h_{ij}+\delta_{ij}(d)$, the pair condition of Lemma~1. Required: both robots use
the same state (A2--A3), apply commands over the same period (A4), the row is enabled for both (the range
test is symmetric), and both QPs are feasible. In simulation the first two hold by construction (one
call on the joint state), so ``decentralised'' refers to how the computation is split. Out-of-range
pairs: a neighbour beyond 1.5\,m needs at least 12 periods at $\le2V_p\Delta\approx0.10$\,m per period to
reach $D_{rr}$, and rows are rebuilt each period. In the final test 0.008\% of robot-pair periods
were beyond 1.5\,m, so limited communication was not tested.
\confirm{the summation argument and its conditions, and the out-of-range argument.}

\section*{C4. Criterion (a) and the error-rate structure}
Usefulness at density $K$ requires (a), (b) and (c) (intersection--union test, so no split across
components). Each component is tested one-sided at $\alpha_K=0.05/3$, a Bonferroni split over the three
density claims, giving familywise error $\le0.05$ over the three usefulness claims.
(a) is decided by the exact one-sided McNemar test: with $n_{10}$ pairs failing off only and $n_{01}$
failing on only, $p=P(X\ge n_{10})$, $X\sim\mathrm{Bin}(n_{10}+n_{01},\tfrac12)$
(\texttt{locked\_stats.mcnemar\_one\_sided\_p}). (b) is decided by the Newcombe bound (C5) against the
5\,pp margin; (c) by the bootstrap bound (C6) against 0.01\,m.
\confirm{the IUT logic, the Bonferroni denominator 3 (densities, not components), and the McNemar direction.}

\section*{C5. Newcombe paired interval: method 8 versus method 10}
With $a,b,c,d$ the paired $2\times2$ counts, $p_1=(a+b)/n$, $p_2=(a+c)/n$, Wilson limits $(l_i,u_i)$ and
\[
\hat\phi=\frac{ad-bc}{\sqrt{(a+b)(c+d)(a+c)(b+d)}},\qquad
\theta-\sqrt{(p_1-l_1)^2-2\hat\phi(p_1-l_1)(u_2-p_2)+(u_2-p_2)^2}
\]
is the lower limit (upper symmetric). This is Newcombe's method 8. Method 10 replaces $ad-bc$ by
$\max(ad-bc-n/2,0)$ when $ad-bc>0$. The locked protocol labels its interval ``method 10'' but describes
it as ``without continuity correction'', i.e.\ method 8; the code
(function \texttt{newcombe\_paired\_ci} in \url{cbf_study/stats.py}, whose docstring also says ``method 10'') implements method 8.
Recomputing every bound with method 10 changes it by $<0.001$\,pp and no criterion
(\url{scripts/newcombe_method10_check.py}). The B-match comparison (Table 4) uses method 10.
When a marginal variance is zero (B-match never fails), $\hat\phi$ is set to 0.
\confirm{the method numbering against Newcombe (1998), the treatment of the inconsistency, and $\hat\phi=0$ in the degenerate case.}

\section*{C6. Criterion (c): formation estimand and bootstrap}
For a pair with executed periods $T_{\rm off},T_{\rm on}$, $L=\min(T_{\rm off},T_{\rm on})$ and
$\Delta e=\frac1L\sum_{t=1}^{L}[e_{\rm on}(t)-e_{\rm off}(t)]$, $e(t)=\tfrac13\sum_r\|x_r(t)-\bar x(t)-s_r\|$;
pairs with $L=0$ are excluded (none occurred). Estimand: $\overline{\Delta e}=\frac19\sum_p\mathrm{mean}_{i\in\mathcal E_p}\Delta e_i$.
Bound: resample pairs with replacement within each actor, recompute the equal-actor mean, $B=20{,}000$,
seed 20261026 (\url{numpy.random.default_rng}), upper bound = $1-0.05/3$ quantile (\texttt{method='linear'});
a normal bound from per-actor variances is a cross-check (function \texttt{formation\_noninferiority} in \url{cbf_study/locked_stats.py}).
Intervals reflect layout sampling for nine fixed actors from three training runs, not variation
across training runs.
\confirm{the estimand, the stratified resampling and the interpretation of the interval.}

\section*{Additional: infeasibility (exploratory, for information)}
At all 433 infeasible robot states, the binding row sat at its stand-off $h^\star=\delta/\alpha$
($|h-h^\star|\le4.3\times10^{-8}$\,m$^2$). An exact vertex enumeration of
$t^\star=\max_{u\in\mathcal U}\min_r(n_r^\top u-c_r)$ gives $t^\star\in[-6.7\times10^{-7},-5.6\times10^{-10}]$; every
minimal conflicting set has three rows (407: two neighbour + one obstacle; 26: one neighbour + two
obstacles), never the command box; 7 cases lie within the solver tolerance $10^{-9}$.
(\texttt{scripts/infeasibility\_conflicts.py}.) A comment on whether this ``three-sided pinch at the
stand-off'' reading is sound would help.

\section*{Sign-off}
\begin{tabular}{>{\raggedright}p{0.36\linewidth}p{0.12\linewidth}p{0.42\linewidth}}
\toprule
Item & Confirmed (Y/N) & Comments \\ \midrule
C1 Lemma 1 and induction & & \\[6pt]
C2 Margins $\delta_o,\delta_{ij}$ & & \\[6pt]
C3 Equal split and ranges & & \\[6pt]
C4 IUT, Bonferroni, McNemar & & \\[6pt]
C5 Newcombe method 8 / 10 & & \\[6pt]
C6 Formation estimand, bootstrap & & \\[6pt]
Margins 5\,pp and 0.01\,m acceptable? & & \\[6pt]
\bottomrule
\end{tabular}

\medskip\noindent Reviewer: \rule{5cm}{0.4pt}\qquad Date: \rule{3cm}{0.4pt}
\end{document}
